\documentclass[a4paper,6pt]{jsarticle}
\usepackage[top=10mm,bottom=10mm,left=10mm,right=10mm]{geometry}
\usepackage{enumitem}
\setlist{nosep}
\pagestyle{empty}

% フォントサイズを最小に調整
\renewcommand{\normalsize}{\fontsize{6.5}{8.5}\selectfont}
\renewcommand{\large}{\fontsize{8}{10}\selectfont}

% インデント幅を最小に調整
\setlist[description]{labelindent=0.2em, leftmargin=1.2em, itemsep=-0.5pt, parsep=0pt, topsep=0pt}

% セクション間のスペースを最小に
\usepackage{titlesec}
\titlespacing*{\section}{0pt}{4pt}{2pt}
\titlespacing*{\subsubsection}{0pt}{2pt}{1pt}

\begin{document}

\vspace*{-10mm}
\begin{flushright}
  \renewcommand{\arraystretch}{3} % 行の高さを2倍に
  \begin{tabular}{|p{2cm}|p{2cm}|p{6cm}|}
    \hline
    \textbf{学科} & \textbf{出席番号} & \textbf{氏名} \\
    \hline
      　　　　　& 　　　　　　& 　　　　　　　　　　　\\
    \hline
  \end{tabular}
\end{flushright}
\vspace{5mm}

\section*{2025年度 情報リテラシー 前期中間試験  問題用紙}

\subsection*{注意事項}

\begin{itemize}
  \item マークシートの学年は、1を記入・マークしてください。
  \item マークシートのクラスは、M:01、E:02、I:03、C:04、A:05の該当する数字を記入・マークしてください。
  \item マークシートの番号は、出席番号を3桁に変換して記入・マークしてください。
  \item 試験時間は \textbf{50分} です。\textbf{持ち込み不可}（教科書・ノート・電子機器など使用できません）。
  \item 問題は全部で \textbf{25問} あります。各問題は \textbf{4点} です。\textbf{100点満点} です。
  \item 解答は\textbf{解答用紙}の解答欄に \textbf{正確に記入}してください。枠外の解答は採点対象にはなりません。
  \item 到達目標1：問1, 問2, 問3, 問4, 問5
\end{itemize}




\subsection*{問1 学内システム活用とアカウント管理}

\subsubsection*{1 学内の複数システムで共通して使用できるアカウントは何か。}
\begin{description}
  \item[A] WebClass
  \item[B] Teams
  \item[C] Forms
  \item[D] 統一アカウント
\end{description}

\subsubsection*{2 アカウントのセキュリティを向上させるために設定する認証方式は何か。}
\begin{description}
  \item[A] QRコード
  \item[B] 多要素認証
  \item[C] 出席管理システム
  \item[D] Webシラバス
\end{description}

\subsubsection*{3 OneDriveの特徴として最も適切なものはどれか。}
\begin{description}
  \item[A] ファイルをローカルPCにのみ保存できる
  \item[B] ファイルの共同編集やクラウド保存が可能
  \item[C] フォームによるアンケートを作成できる
  \item[D] メール送信やビデオ会議ができる
\end{description}

\subsubsection*{4 クラスの相談や発表準備のグループ会議など、双方向のやり取りに最も適したアプリはどれか。}
\begin{description}
  \item[A] OneDrive
  \item[B] Forms 
  \item[C] Teams  
  \item[D] Outlook
\end{description}

\subsubsection*{5 授業の感想やアンケートを集計したいとき、最も適したアプリはどれか。}
\begin{description}
  \item[A] OneDrive
  \item[B] Forms
  \item[C] Teams
  \item[D] PowerPoint
\end{description}

\subsection*{問2 情報の特性と情報モラル}

\subsubsection*{1 以下のうち、「複製性」に該当する例として最も適切なのはどれか。}
\begin{description}
  \item[A] 写真をネットに投稿したら、別サイトに転載された
  \item[B] 画像の著作権を確認せずに印刷した
  \item[C] 個人情報が書かれたファイルをUSBで持ち出した
  \item[D] メールの宛先を間違えて送信してしまった
\end{description}

\subsubsection*{2 次のうち、「特許権」で保護される対象として最も適切なものはどれか。}
\begin{description}
  \item[A] ロゴや商品パッケージのデザイン
  \item[B] 小説や楽曲などの創作物
  \item[C] 新しい技術や発明のアイデア
  \item[D] サービス名や会社名などの名称
\end{description}

\subsubsection*{3 以下のうち、「意匠権」で保護される対象はどれか。}
\begin{description}
  \item[A] スマートフォンの内部構造のアルゴリズム
  \item[B] パソコンに使われるOSのソースコード
  \item[C] ユーザーインターフェースや機器の見た目のデザイン
  \item[D] クラウドサービスのサーバ構成図
\end{description}

\subsubsection*{4 次のうち、「フィッシング詐欺の例」として最も適切なものはどれか。}
\begin{description}
  \item[A] コンビニでプリペイドカードの番号を電話で伝えさせる詐欺
  \item[B] 警察を名乗って金を振り込ませる電話詐欺
  \item[C] 有名企業のロゴを使った偽メールでパスワードを入力させる行為
  \item[D] スマホのアプリに課金を続けてしまう詐欺的な仕組み
\end{description}

\subsubsection*{5 「オプトイン」の説明として正しいものは何か。}
\begin{description}
  \item[A] 情報が広がりやすい性質
  \item[B] 事前に同意を得ること
  \item[C] 情報が残り続ける性質
  \item[D] 著作権を侵害すること
\end{description}

\newpage
\vspace{-5mm}
\subsection*{問3 情報システムとクラウドコンピューティング}

\subsubsection*{1 次のうち、情報システムの「処理」に該当する操作はどれか。}
\begin{description}
  \item[A] キーボードで氏名を入力する
  \item[B] プリンターで成績表を印刷する
  \item[C] 入力された得点をもとに平均点を計算する
  \item[D] 成績データをハードディスクに保存する
\end{description}

\subsubsection*{2 次のうち、ATMの「出力」に該当するものはどれか。}
\begin{description}
  \item[A] キャッシュカードを挿入する
  \item[B] 預金残高が画面に表示される
  \item[C] 現金を引き出す金額をテンキーで入力する
  \item[D] 引き出し額の指定に対して本人確認を行う
\end{description}

\subsubsection*{3 次のうち、「PaaS」の説明として最も適切なものはどれか。}
\begin{description}
  \item[A] 仮想マシンやネットワークなどインフラを提供するサービス
  \item[B] 完成されたアプリケーションをそのまま使えるサービス
  \item[C] アプリ開発に必要な環境（OSやデータベース）を提供するサービス
  \item[D] ストレージやデバイスを個人で所有せずに使うレンタルサービス
\end{description}

\subsubsection*{4 必要に応じてシステム規模を簡単に増減できるクラウドの特徴は何か。}
\begin{description}
  \item[A] 処理性
  \item[B] スケーラビリティ
  \item[C] 保存性
  \item[D] 出力性
\end{description}

\subsubsection*{5 次のうち、「SaaS」の活用例として最も適切なものはどれか。}
\begin{description}
  \item[A] AWS EC2で仮想サーバをレンタルする
  \item[B] Google App Engineでアプリを開発・公開する
  \item[C] 自社でサーバを購入してシステムを構築する
  \item[D] ブラウザ上でGmailを使ってメールを送受信する
\end{description}

\vspace{-2mm}
\subsection*{問4 PC設定とアカウント管理}

\subsubsection*{1 Windowsのローカルアカウントについて正しい説明はどれか。}
\begin{description}
  \item[A] Microsoftアカウントが必須のオンライン専用アカウント
  \item[B] ネットワーク接続がなくても使用できる、PC内のアカウント
  \item[C] 自動的に削除される一時アカウント
  \item[D] クラウド上にデータを自動同期するアカウント
\end{description}

\subsubsection*{2 Windowsの初期設定でコマンドプロンプトを表示する操作として正しいものはどれか。}
\begin{description}
  \item[A] Ctrl + Alt + Delete を押す
  \item[B] Shift + F10 を押す
  \item[C] Windows + R を押す
  \item[D] 先生を呼ぶ
\end{description}

\subsubsection*{3 「oobe\textbackslash{}BypassNRO」コマンドの目的として最も適切なものはどれか。}
\begin{description}
  \item[A] パスワード入力を省略するための設定
  \item[B] ローカルアカウント作成時にネットワーク接続を回避する
  \item[C] セキュリティの質問をスキップするための機能
  \item[D] 管理者権限を取得するためのコマンド
\end{description}

\subsubsection*{4 ネットワーク接続の詳細設定画面を開くためのコマンドとして正しいものはどれか。}
\begin{description}
  \item[A] netplwiz
  \item[B] msconfig
  \item[C] ncpa.cpl
  \item[D] ls
\end{description}

\subsubsection*{5 ローカルアカウントを作成し、ネットワークを元に戻すまでの手順を正しい順序で並べよ。（記入例：A B C D E）}

（　　） → （　　） → （　　） → （　　） → （　　）

\noindent 選択肢： 
\begin{description}
  \item[A] ネットワークを有効に戻す  
  \item[B] セキュリティの質問を設定する  
  \item[C] ネットワークを無効にする  
  \item[D] コマンド「oobe\textbackslash{}BypassNRO」を実行する  
  \item[E] ユーザー名を設定する  
\end{description}

\vspace{-2mm}
\subsection*{問5 Officeアプリケーションとクラウドサービス}

\subsubsection*{1 自分のPCにOfficeを新規インストールする際、まず行うべき適切な操作はどれか。}
\begin{description}
  \item[A] Wordアプリを開き、ローカルアカウントを作成する
  \item[B] portal.office.com にアクセスして、Microsoftアカウントでサインインする
  \item[C] インストーラーをUSBで別のPCからコピーして実行する
  \item[D] Microsoftの株を買う
\end{description}

\subsubsection*{2 Officeのインストール時に、インストーラーが「ネットワークに接続できません」と表示された。最初に確認すべきこととして最も適切なのはどれか。}
\begin{description}
  \item[A] Officeが最新バージョンであるかを確認する
  \item[B] 有線LANなどでPCがインターネットに接続されているか確認する
  \item[C] Wordのショートカットが正しくデスクトップに作成されているか確認する
  \item[D] Windowsが最新の状態にアップデートされているか確認する
\end{description}

\subsubsection*{3 Wordの標準的なファイル形式はどれか。}
\begin{description}
  \item[A] .txt
  \item[B] .docx
  \item[C] .xlsx
  \item[D] .pptx
\end{description}

\subsubsection*{4 ファイルをクラウド上に保存・共有できるMicrosoft製ストレージサービスは何か。}
\begin{description}
  \item[A] iCloud Drive
  \item[B] DropBox
  \item[C] OneDrive
  \item[D] Google Drive
\end{description}

\subsubsection*{5 Officeインストールの手順を正しい順序で並べよ。（記入例：A B C D）}
（　　） → （　　） → （　　） → （　　）

\noindent 選択肢： 
\begin{description}
  \item[A] インストーラー実行
  \item[B] portal.office.com にアクセス
  \item[C] 有線LAN接続
  \item[D] Microsoftアカウントでサインイン
\end{description}

\end{document}